\section{v1.17：把多块更新变成可恢复事务}

\begin{keyidea}
程序顺序只约束 CPU 怎样发出写入，不能自动约束断电后哪些字节仍在介质上。
ordered metadata journal 用 ``data stable''、``prepared log stable'' 和
``commit stable'' 三条持久顺序，把一次多块元数据修改收敛成旧状态或完整
新状态；checkpoint 可以重复执行，因而恢复本身再次断电也不会制造第三种状态。
\end{keyidea}

\subsection{从事后检查到写前日志}

早期 Unix 文件系统直接覆盖目录、inode 与 free list。异常断电后，\code{fsck}
只能遍历整棵命名空间，重新推导 ``哪些对象可达、哪些块被重复使用''。这种
方法能修补结构，却不知道半完成的 rename 原本应该选择旧名字还是新名字；
介质越大，完整扫描也越慢。

写前日志（write-ahead logging，WAL）改变了证据顺序：在覆盖正常位置
（home block）前，先把本事务的完整新元数据写到独立区域；只有这些材料已经
稳定，才写一条 commit record。恢复程序不猜测业务意图，只判断 commit：

\begin{itemize}
  \item 没有完整 commit：home metadata 的旧状态仍然有效，丢弃 prepared log；
  \item 有完整 commit：journal payload 是权威新状态，把它复制到 home；
  \item 记录字段、范围或校验不合法：拒绝 mount，不能把损坏伪装成某个事务。
\end{itemize}

数据库与文件系统都使用 WAL，但对象不同。数据库日志常记录页修改、逻辑操作
或 undo/redo；本项目记录 4\,KiB 元数据块的完整新镜像，只需要 redo。

\subsection{为什么选择 ordered metadata，而不是记录全部数据}

常见 journal 策略可分成三类：

\begin{description}
  \item[data journal] 普通文件数据和元数据都入日志，语义强，写放大和日志
    容量也最大；
  \item[ordered metadata] 只把元数据放入日志，但 metadata commit 前必须
    先稳定与之相关的数据；
  \item[writeback metadata] 元数据仍有结构原子性，却不保证数据先于引用它
    的 inode/目录落盘。
\end{description}

v1.16 已经建立 dirty/writeback/error page、BlockRequest 和 ATA FLUSH
边界。v1.17 选择 ordered mode，是因为它能直接复用这条 ``相关数据已稳定''
证据，同时把教学范围限制在元数据事务。它保证目录、inode、位图和指针不会
形成半事务，不承诺普通数据旧内容可回滚。

例如覆盖一个已经分配的数据块后，在 metadata commit 前断电：旧 inode 仍然
指向合法块，文件系统结构一致，但块内容可能已经是新字节。这不是 replay
错误，而是 metadata-only 的明确边界。

\subsection{格式 3 的固定区域}

rootfs 仍使用 256\,MiB 分区和 512\,B 文件系统块。盘面 magic 从格式 2
升级为 \code{OSRFV003}，并固定保留 256 个 journal block：

\begin{center}
\begin{tabular}{lll}
\toprule
相对块 & 数量 & 含义\\
\midrule
0 & 1 & superblock home\\
1--4226 & 4226 & 位图与 inode table\\
4227--4482 & 256 & journal\\
4483--524287 & 519805 & 目录和普通文件数据\\
\bottomrule
\end{tabular}
\end{center}

固定位置很重要：mount 必须在读取并信任 home superblock 之前恢复，而恢复
不能反过来依赖一个可能尚未 checkpoint 的 superblock 字段来寻找 journal。

journal 内部布局为：

\begin{center}
\begin{tabular}{lll}
\toprule
journal 相对块 & 数量 & 内容\\
\midrule
0 & 1 & header\\
1--4 & 4 & descriptor blocks\\
5--128 & 124 & metadata payload blocks\\
129 & 1 & commit record\\
130--255 & 126 & 冻结保留区\\
\bottomrule
\end{tabular}
\end{center}

一个 descriptor entry 是 16 字节：64 位 target relative block、32 位
payload CRC32 与 32 位保留零。descriptor 最后 4 字节保存自身 CRC，因此
每块容量是
\[
  \left\lfloor\frac{512-4}{16}\right\rfloor=31,
\]
四块公开常量给出的总 credit 是
\[
  4\times31=124.
\]
其余 descriptor 字节必须为零并纳入 CRC。这里的 124 来自冻结格式，而不是
随手散落在代码中的数字；格式工具、Kernel 和测试读取同一显式常量契约。

\begin{contractbox}
教材中的容量必须以实现的冻结字段为准。若未来把每个 descriptor 块扩到理论
最大 entry 数，那是盘面格式升级，不能只改一个循环上限。
\end{contractbox}

\subsection{header、descriptor、payload 与 commit 怎样互相约束}

header 保存 magic、格式版本、64 位 sequence、entry count、descriptor count
和 CRC32。descriptor 给每个 payload 指定 home target 与 payload CRC。commit
再次保存独立 magic、版本、sequence、entry count、header CRC 和自身 CRC。

恢复只有在以下检查全部通过后才能写第一个 home block：

\begin{enumerate}
  \item header 与 commit 的 magic/version 正确；
  \item sequence 非零且两处相同；
  \item entry count 在 1 到 124，且两处相同；
  \item descriptor count 等于冻结值；
  \item 每个 descriptor CRC 与 payload CRC 正确；
  \item target 在 rootfs 内、不落入 journal、同一事务内不重复；
  \item reserved 字段与未使用区域满足格式约束；
  \item commit 的 header CRC 与 commit CRC 正确。
\end{enumerate}

先验证全部记录、再统一 checkpoint，避免恢复前半段后才在后半段发现损坏。
CRC32 用于检测随机位翻转、截断与错块，不是密码学认证；面对主动构造碰撞的
攻击者，边界检查仍负责保证只能拒绝，不能越界或 panic。

\subsection{credit 是空间预留，也是失败原子性}

一个操作如果修改到一半才发现 journal 放不下，Kernel 需要撤销所有内存状态，
且相关普通数据可能已经变化。transaction credit 把容量判断前移：

\begin{itemize}
  \item Begin 只接受 1--124；
  \item 首次 stage 一个新 target 消耗一个 credit；
  \item 再次 stage 相同 target 只替换 overlay，不重复计费；
  \item 超过 reservation 返回 CreditsExhausted，home metadata 仍未改变。
\end{itemize}

当前 rootfs 串行一个 mutation，并预留固定上限。Begin 同时保存 superblock、
统计与 allocation hint 快照。commit 前的逻辑、credit 或设备错误会丢弃
overlay 并恢复快照。若未来支持多个并发事务，应先冻结锁顺序、reservation
所有权和 checkpoint 调度，不能简单把固定数组复制到每个 Thread。

\subsection{overlay 让事务读到自己的写}

Create 可能先更新 inode block，随后又因目录插入读取同一块。如果 read 只看
BlockCache，就会观察旧字节。v1.17 的 metadata read 顺序为：

\begin{center}
\begin{tikzpicture}[node distance=9mm and 13mm]
  \node[os state] (read) {metadata read};
  \node[os block,right=of read] (overlay) {journal overlay};
  \node[os block,right=of overlay] (cache) {BlockCache/device};
  \node[os state,below=of overlay] (value) {完整 4\,KiB block};
  \draw[os arrow] (read) -- (overlay);
  \draw[os arrow] (overlay) -- node[above]{miss} (cache);
  \draw[os arrow] (overlay) -- node[left]{hit} (value);
  \draw[os arrow] (cache) -- (value);
\end{tikzpicture}
\end{center}
\diagramnote{overlay 是固定容量 512\,B 整块镜像，不依赖动态分配，也不保存容易错序的小 patch。}

\subsection{四次 FLUSH 建立的持久 happens-before}

成功事务严格执行：

\begin{enumerate}
  \item 把相关普通数据写到 home，执行 FLUSH；
  \item 写 header、descriptor 和 payload，执行 prepared FLUSH；
  \item 写 commit record，执行 commit FLUSH；
  \item payload checkpoint 到 metadata home，执行 checkpoint FLUSH；
  \item 清零 header/commit，执行 clear FLUSH。
\end{enumerate}

从性能角度看 FLUSH 很贵，但这里的每一个都有语义：

\[
\text{related data stable}
\prec
\text{metadata commit stable}
\]
\[
\text{prepared payload stable}
\prec
\text{commit stable}
\]
\[
\text{home checkpoint stable}
\prec
\text{journal reusable}
\]

CPU 按顺序调用 Write 并不能替代这些关系。ATA/QEMU cache 可以让命令完成早于
掉电持久化；没有 FLUSH 的 ``先写 payload、后写 commit'' 只是提交顺序。

\subsection{把所有断电点归约为两种状态}

\begin{center}
\begin{tabularx}{\textwidth}{p{0.24\textwidth}XX}
\toprule
断电窗口 & mount 看到的证据 & 恢复结果\\
\midrule
data 或 prepared 写入前/中 &
没有有效 commit &
保留旧 metadata，丢弃完整 prepared 或拒绝损坏格式\\
prepared FLUSH 后 &
payload 完整但 commit 缺失 &
清理 incomplete，保留旧状态\\
commit 写入中 &
字段或 CRC 不完整 &
不能伪造 commit；按严格分类丢弃或拒绝\\
commit FLUSH 后 &
header、payload、commit 全部匹配 &
replay 全部 target，得到完整新状态\\
checkpoint 中 &
部分 home 新、commit 仍完整 &
重新 replay 全部 target，得到完整新状态\\
checkpoint 后、clear 前 &
home 已全新，commit 仍完整 &
再次写同样字节，结果不变\\
\bottomrule
\end{tabularx}
\end{center}

这里的线性化点是 commit FLUSH 成功。它之前恢复选择旧 metadata，它之后恢复
必须选择新 metadata。

\subsection{幂等 replay 让恢复也可以被中断}

checkpoint 不是不可中断的魔法操作。若恢复写完第 37 个 target 后再次断电，
下一次 mount 会从第一个 target 重新写。本项目保存完整新块，因此：

\[
\operatorname{Apply}(\operatorname{Apply}(H,P),P)
=\operatorname{Apply}(H,P).
\]

checkpoint 或其 FLUSH 失败时不能清除 commit；当前文件系统实例被冻结，下一
次 mount 仍有权威恢复材料。只有 home 全部稳定后才能清 journal，并再次
FLUSH 清理结果。

\subsection{生产源码的责任分界}

\code{fs/root\_journal.*} 只认识固定块设备、盘面记录、overlay、CRC、FLUSH、
checkpoint 和 replay，不认识路径、inode 语义或 Process。\code{fs/root\_file\_system.*}
决定哪些写是 metadata、何时先同步数据、如何保存事务快照以及设备失败后何时
冻结。\code{fs/root\_file\_system\_format.*} 只编码/解析小端字段。

所有盘面整数通过显式 Store/Load 写入字节数组，不能把 C++ struct 直接
\code{reinterpret\_cast} 到磁盘。后者会把 padding、alignment、宿主端序与
编译器 ABI 偷偷变成文件格式。

Initialize 的顺序同样是协议：

\begin{enumerate}
  \item 绑定 block device 与固定 journal 位置；
  \item Recover journal；
  \item 读取并验证 home superblock；
  \item 初始化 cache 与运行期状态；
  \item 输出 journal ready 与有界恢复摘要。
\end{enumerate}

\subsection{独立 fsck 为什么仍不可删除}

journal 证明最后一个事务的旧/新选择；fsck 证明最终全局结构。独立 Python
fsck 从根目录重算可达 inode、data block、link 与 bitmap，不调用 Kernel
RootJournal parser。两个 oracle 若共享实现，同一个 target 换算错误可能让
生产代码和测试一起说 ``正确''。

\subsection{1000 个断电点和十万步随机模型}

故障块设备给每一次 Write/Flush 编号，在指定编号完成后模拟断电。集成测试从
同一 old image 重建 1000 次，恢复后逐 target 验证：

\[
(\forall i,\ H_i=Old_i)
\oplus
(\forall i,\ H_i=New_i).
\]

不允许部分旧、部分新。恢复后的镜像再执行一次 Recover，必须报告 Clean。

随机模型用固定 seed 执行 100000 步 Begin、Stage、Abort、Commit 和非法
操作。独立 oracle 跟踪 active、staged 与 durable；durable 只允许在 commit
后变化，失败操作不能发布状态。固定 seed 使任何失败可复现，而不是用 ``随机
跑过一次'' 代替证明。

\subsection{日志只记录低频语义边界}

逐 payload、逐 CRC 或逐 ATA word 打印会经串口改变 PIT/IRQ 时序。生产日志
只输出 journal ready、credit capacity、commit、replay、discard 与 checksum
failure 汇总，并带来宾单调时间。宿主 QEMU 的到达时间帮助定位停滞，不充当
来宾事件顺序。

\subsection{本阶段没有偷偷承诺什么}

v1.17 不包含 data journal、后台 checkpoint、并行 transaction、SMP、
\code{fsync}/\code{fdatasync}/\code{msync}、snapshot、在线扩容、AHCI 或
NVMe。它也不声称 CRC 能抵抗主动篡改。

这些边界让结论保持精确：现在已经能证明多块 metadata mutation 在覆盖的任意
断电点后只有旧/新两种结构状态；下一阶段应冻结 ABI、补齐观察面和畸形输入
审计，而不是继续往 journal 里塞新的文件系统机制。
